%----------------------------------------------------------------------
% 英文摘要
%----------------------------------------------------------------------

% 把英文摘要寫在裡面
\begin{enAbstract}
Two-dimensional transition-metal dichalcogenides (TMDs) are promising materials for optoelectronic applications due to their direct band gaps and strong excitonic effects. Through the electron–hole exchange interaction, the valley-mixed excitonic branches possess different transition dipole orientations in the center-of-mass (COM) momentum space. Therefore, utilizing structured optical fields to excite these excitonic states has become an important research direction. Vector vortex beams (VVBs), formed by the superposition of two Laguerre–Gaussian beams, provide a versatile platform for tailoring structured light. Through spin–orbit interaction in the optical field, the longitudinal component of VVBs exhibits characteristic angular geometry, which corresponds to exciton transition-rate distributions in COM momentum space~\cite{sanchez2024enhanced}. Motivated by this connection, this thesis investigates generalized VVB configurations and their interaction with excitons.

For VVBs composed of modes with equal orbital angular momentum (OAM) magnitudes, the total angular momentum (TAM) difference and OAM difference are necessarily even integers, resulting in symmetric exciton transition-rate distributions. By extending to VVBs with unequal OAM magnitudes, we demonstrate that the TAM difference and OAM difference independently determine the angular geometry of the longitudinal and transverse components, respectively, manifested as intensity and polarization distributions. This governing mechanism remains valid for VVBs with unequal OAM magnitudes as long as sufficient spatial overlap between the constituent beams is maintained. In particular, odd TAM or OAM differences enable directional-selective excitation of finite-momentum excitons.

Finally, we construct exciton wave packets while retaining the internal exciton wave function. We find that, for valley-mixed excitons, the phase contribution originating from the internal exciton structure becomes a relative phase between different valley components and plays a crucial role in determining the exciton wave-packet structure. Therefore, further numerical calculations are required for a complete description of the wave-packet structure.
    \enAbsKeywords
\end{enAbstract}





